% =============================================================================
\section{Modelo Matemático}
% =============================================================================

% ---------------------------------------------------------------------------
\subsection{Clasificación del sistema de colas}
% ---------------------------------------------------------------------------

El sistema de La Cocina de Kojo puede clasificarse, usando la notación de
Kendall, como un sistema \textbf{M/G/$c$}:
\begin{itemize}
  \item \textbf{M} (Markoviano): las llegadas de clientes forman un proceso
    de Poisson, cuyos tiempos entre eventos son variables exponenciales
    independientes --- memoria nula entre llegadas sucesivas.
  \item \textbf{G} (General): los tiempos de servicio siguen distribuciones
    uniformes, no exponenciales.
  \item \textbf{$c$}: hay $c \in \{2, 3\}$ servidores paralelos con cola
    única FIFO.
\end{itemize}

Dado que la tasa de llegadas no es constante a lo largo del día, el modelo
global es en realidad un sistema M/G/$c$ \emph{no estacionario}.
El análisis teórico que sigue trata cada período (pico / fuera de pico) de
forma independiente asumiendo condiciones de régimen estacionario dentro de
cada tramo.

% ---------------------------------------------------------------------------
\subsection{Parámetros del servicio}
% ---------------------------------------------------------------------------

\paragraph{Tiempo de servicio esperado.}
Con probabilidad $p = 0{,}5$ el cliente pide sándwich
($S_{\text{sand}} \sim \mathcal{U}[3,5]$) y con probabilidad $1-p = 0{,}5$
pide sushi ($S_{\text{sush}} \sim \mathcal{U}[5,8]$).
La esperanza del tiempo de servicio es:
\begin{equation}
  \EE[S] = p\,\EE[S_{\text{sand}}] + (1-p)\,\EE[S_{\text{sush}}]
         = 0{,}5 \cdot \frac{3+5}{2} + 0{,}5 \cdot \frac{5+8}{2}
         = 0{,}5 \times 4 + 0{,}5 \times 6{,}5
         = \mathbf{5{,}25} \text{ min.}
  \label{eq:ES}
\end{equation}

La tasa de servicio de un empleado es:
\[
  \mu = \frac{1}{\EE[S]} = \frac{1}{5{,}25} \approx 0{,}1905 \text{ clientes/min.}
\]

\paragraph{Segundo momento y varianza del servicio.}
Para $X \sim \mathcal{U}[a,b]$ se tiene
$\EE[X^2] = (a^2 + ab + b^2)/3$.
Por tanto:
\[
  \EE[S_{\text{sand}}^2] = \frac{9 + 15 + 25}{3} = \frac{49}{3} \approx 16{,}33, \qquad
  \EE[S_{\text{sush}}^2] = \frac{25 + 40 + 64}{3} = 43.
\]
\[
  \EE[S^2] = 0{,}5 \times \frac{49}{3} + 0{,}5 \times 43
           = 8{,}167 + 21{,}5 = \mathbf{29{,}67} \text{ min}^2.
\]
\[
  \VV[S] = \EE[S^2] - (\EE[S])^2 = 29{,}67 - 5{,}25^2 = 29{,}67 - 27{,}5625
         \approx \mathbf{2{,}11} \text{ min}^2, \quad
  \sigma_S \approx 1{,}45 \text{ min.}
\]

% ---------------------------------------------------------------------------
\subsection{Utilización teórica y estabilidad}
% ---------------------------------------------------------------------------

La \textbf{utilización} de cada servidor en el período analizado es:
\begin{equation}
  \rho = \frac{\lambda}{c \cdot \mu},
  \label{eq:rho}
\end{equation}
con la condición de estabilidad $\rho < 1$.
La \cref{tab:utilizacion_teorica} muestra los valores para cada combinación
de período y escenario.

\begin{table}[H]
  \centering
  \caption{Utilización teórica $\rho$ por período y escenario.}
  \label{tab:utilizacion_teorica}
  \begin{tabular}{llcccc}
    \toprule
    \textbf{Período} & \textbf{Escenario} & $\lambda$ & $c$ & $\mu$ & $\rho = \lambda/(c\mu)$ \\
    \midrule
    Fuera de pico & A & 0{,}17 & 2 & 0{,}1905 & 0{,}446 \\
    Fuera de pico & B & 0{,}17 & 2 & 0{,}1905 & 0{,}446 \\
    Pico          & A & 0{,}30 & 2 & 0{,}1905 & 0{,}787 \\
    Pico          & B & 0{,}30 & 3 & 0{,}1905 & 0{,}525 \\
    \bottomrule
  \end{tabular}
\end{table}

En todos los casos $\rho < 1$: el sistema es estable en cada período.
La carga más exigente corresponde al período pico del Escenario A ($\rho =
0{,}787$), lo que explica la acumulación de cola observada en la simulación.

% ---------------------------------------------------------------------------
\subsection{Aproximación M/M/c como cota teórica}
% ---------------------------------------------------------------------------

Para obtener predicciones analíticas se aproxima el sistema M/G/$c$ por un
sistema M/M/$c$ con la misma tasa media de servicio.
Aunque los tiempos de servicio no son exponenciales, este modelo proporciona
una cota de referencia útil.

\paragraph{Fórmula de Erlang~C.}
La probabilidad de que un cliente deba esperar en cola es~\citep{banks2010discrete}:
\begin{equation}
  C(c,\,a) = \frac{\displaystyle\frac{a^c}{c!}\cdot\frac{c}{c - a}}
                  {\displaystyle\sum_{n=0}^{c-1}\frac{a^n}{n!}
                   + \frac{a^c}{c!}\cdot\frac{c}{c - a}},
  \label{eq:erlangC}
\end{equation}
donde $a = \lambda/\mu$ es la intensidad de tráfico.

\paragraph{Tiempo de espera en cola (M/M/$c$).}
\begin{equation}
  W_q^{\text{MM}c} = \frac{C(c,\,a)}{c\,\mu\,(1-\rho)}.
  \label{eq:Wq_mmc}
\end{equation}

\paragraph{Longitud media de cola.}
Por la Ley de Little~\citep{little1961proof}:
\begin{equation}
  L_q = \lambda\,W_q.
  \label{eq:little}
\end{equation}

Los resultados se muestran en la \cref{tab:teoria_vs_sim}.
Los valores teóricos se calculan como media ponderada por tiempo de los dos períodos
del día: $W_q^{\text{día}} = \frac{240}{660}\,W_q^{\text{pico}} + \frac{420}{660}\,W_q^{\text{fuera}}$,
con $\lambda_{\text{eff}} = \frac{240 \times 0{,}30 + 420 \times 0{,}17}{660} = 0{,}217$ clientes/min.

\begin{table}[H]
  \centering
  \caption{Comparación teoría M/M/$c$ vs.\ simulación ---
           media ponderada del día completo.}
  \label{tab:teoria_vs_sim}
  \begin{tabular}{llccc}
    \toprule
    \textbf{Escenario} & \textbf{Indicador}
      & \textbf{M/M/$c$ teórico}
      & \textbf{Simulación (media)}
      & \textbf{Sim $-$ Teo (\%)} \\
    \midrule
    A ($c_{\text{pico}}=2$) & $W_q$ (min)      & 3{,}95 & 2{,}19 & $-$44{,}6 \\
    A ($c_{\text{pico}}=2$) & $L_q$ (clientes) & 0{,}86 & 0{,}48 & $-$44{,}2 \\
    B ($c_{\text{pico}}=3$) & $W_q$ (min)      & 1{,}18 & 0{,}66 & $-$44{,}1 \\
    B ($c_{\text{pico}}=3$) & $L_q$ (clientes) & 0{,}26 & 0{,}15 & $-$42{,}3 \\
    \bottomrule
  \end{tabular}
\end{table}

% ---------------------------------------------------------------------------
\subsection{Supuestos, restricciones y fuentes de discrepancia}
% ---------------------------------------------------------------------------

Las diferencias entre el modelo M/M/$c$ y la simulación se explican por
varios factores:

\begin{enumerate}
  \item \textbf{Distribución del servicio (factor dominante).}
    Los tiempos de servicio son uniformes, no exponenciales.
    La distribución exponencial tiene coeficiente de variación $C_v = 1$,
    mientras que $\mathcal{U}[3,5]$ tiene $C_v = 0{,}14$,
    $\mathcal{U}[5,8]$ tiene $C_v = 0{,}13$, y la mezcla combinada tiene
    $C_v = \sigma_S / \EE[S] = 1{,}45 / 5{,}25 \approx 0{,}28$.
    Tiempos de servicio mucho menos variables que los exponenciales generan
    colas notablemente menores, como predice la fórmula de
    Pollaczek--Khinchine~\citep{ross2013simulation}.
    Este factor explica que la simulación dé resultados aprox.\ \SI{44}{\percent}
    menores que M/M/$c$ de manera consistente en ambos escenarios.

  \item \textbf{No estacionariedad.}
    La tasa $\lambda(t)$ cambia a lo largo del día.
    El modelo estacionario sobreestima el efecto del pico porque asume que la
    tasa alta ha estado vigente indefinidamente, acumulando más cola de la real.

  \item \textbf{Condición transitoria.}
    Al inicio de cada período pico el sistema parte de una cola residual, no del
    estado de régimen estacionario que asume el modelo analítico.

  \item \textbf{Jornada finita.}
    La Ley de Little y la fórmula de Erlang C son resultados de largo plazo
    (régimen estacionario); aquí el horizonte es de 660 min.
\end{enumerate}

\paragraph{Verificación con la Ley de Little.}
Se puede verificar internamente la consistencia de la simulación usando
$L_q = \lambda_{\text{eff}}\,W_q$, donde $\lambda_{\text{eff}}$ es la tasa
efectiva de clientes atendidos.
Para el Escenario A: $0{,}482 \approx (144{,}2 / 660) \times 2{,}185 = 0{,}218 \times 2{,}185 = 0{,}477$ $\checkmark$
